\setcounter{section}{0} 

\section{Structured Spaces}

This semester, you will be introduced to certain structures that can be given to a simple set (which we will often call a \qt{space} $X$ in this context). You can think of this as a hierarchy, where each level grants us more powerful mathematical tools:

\begin{enumerate}
    \item \textbf{Topological spaces:} \textcolor{eGray}{(Focus in the 4th semester)}
    \item \textbf{Metric spaces $(X, d)$:} Defines a mere \emph{distance} between points. Important: $X$ does not necessarily have to be a vector space here! \textcolor{eGray}{(Focus in Analysis)}
    \item \textbf{Normed vector spaces $(X, \norm{\cdot})$:} Defines the \emph{length} of a vector. This presupposes a linear structure (i.e., $X$ must be a vector space). \textcolor{eGray}{(Focus in LinAlg / Analysis)}
    \item \textbf{Inner product spaces $(X, \langle \cdot, \cdot \rangle)$:} Defines both lengths and \emph{angles}! This is the richest structure. \textcolor{eGray}{(Focus in LinAlg)}
\end{enumerate}

\gemininote{Why is this hierarchy important? Because the more powerful structures automatically induce the weaker ones! An inner product always induces a norm via the formula $\norm{x} := \sqrt{\langle x, x \rangle}$. And a norm always induces a metric via the formula $d(x,y) := \norm{x - y}$. The reverse, however, is not generally true!}

\hrulefill

\section{Metric Spaces}

A \textbf{metric space} $(X, d)$ is any set $X$ equipped with a \textbf{metric}. You can think of it as a universal ruler.

Formally, a metric is a function:
\[ d: X \times X \to \mathbb{R}_{\ge 0} \]
such that for all points $x, y, z \in X$, the following three intuitive properties hold:

\begin{enumerate}
    \item \textbf{Definiteness:} $d(x,y) \ge 0$, and $d(x,y) = 0 \iff x=y$. 
    \\ \textcolor{eGray}{In words: \qt{The distance is always positive, and a point has zero distance only to itself.}} 
    
    \item \textbf{Symmetry:} $d(x,y) = d(y,x)$.
    \\ \textcolor{eGray}{In words: \qt{The path from you to me is exactly as long as the path from me to you.}} 
    
    \item \textbf{Triangle Inequality:} $d(x,z) \le d(x,y) + d(y,z)$.
    \\ \textcolor{eGray}{In words: \qt{Taking a detour through a third point can never shorten your trip.}} 
\end{enumerate}

\begin{center}
\begin{tikzpicture}[scale=1.5, very thick]
    \coordinate (X) at (0,0);
    \coordinate (Y) at (2,1.5);
    \coordinate (Z) at (4,0);

    \draw[eMediumBlue, ->] (X) -- node[above left] {$d(x,y)$} (Y);
    \draw[eMediumBlue, ->] (Y) -- node[above right] {$d(y,z)$} (Z);
    \draw[eDarkRed, dashed, ->] (X) -- node[below] {$d(x,z) \le d(x,y) + d(y,z)$} (Z);

    \fill[eSaddleBrown] (X) circle (2pt) node[left] {$x$};
    \fill[eSaddleBrown] (Y) circle (2pt) node[above] {$y$};
    \fill[eSaddleBrown] (Z) circle (2pt) node[right] {$z$};
\end{tikzpicture}
\end{center}

\subsection*{Examples of Metric Spaces}

\begin{example*}[The Discrete Metric]
Let $X$ be any set. The discrete metric is the \qt{digital} approach to measuring distance:
\[
d(x,y) = 
\begin{cases} 
1 & \text{if } x \ne y \\ 
0 & \text{if } x = y 
\end{cases}
\]
\gemininote{This metric is an extremely important tool for constructing counterexamples! It perfectly demonstrates that a metric does not have to match our intuitive geometric understanding of space.}
\end{example*}

\begin{example*}[Metrics in $\R^n$]
We have many choices for measuring distance in $\R^n$. Some common ones are:
\begin{itemize}
    \item \textbf{Manhattan metric ($d_1$):} $d_1(x,y) = \sum_{k=1}^n \abs{x_k - y_k}$ \textcolor{eGray}{(Like a taxi navigating a city grid)}
    \item \textbf{Standard/Euclidean metric ($d_2$):} $d_2(x,y) = \sqrt{\sum_{k=1}^n (x_k - y_k)^2}$ \textcolor{eGray}{(The classical straight-line distance)}
    \item \textcolor{eDarkGreen}{\qt{Supremum metric}} \textbf{($d_3$ or $d_\infty$):} $d_3(x,y) = \sup_{k=1,\dots,n} \abs{x_k - y_k}$ \textcolor{eGray}{(Only the single largest coordinate difference matters)}
\end{itemize}

As a generalization: $d_m(x,y) = \sqrt[m]{\sum_{k=1}^n \abs{x_k - y_k}^m}$ is a valid metric on $\R^n$ for any $m \ge 1$.

\begin{center}
\textbf{Visualizing the \qt{Unit Balls} in $\R^2$:} \\
\vspace{0.3cm}
\begin{tikzpicture}[scale=2]
    \draw[->, thick, eBlack] (-1.5,0) -- (1.5,0) node[right] {$x_1$};
    \draw[->, thick, eBlack] (0,-1.5) -- (0,1.5) node[above] {$x_2$};
    
    \draw[eDarkGreen, very thick, dashed] (-1,-1) rectangle (1,1);
    \node[eDarkGreen] at (1.2, 1.2) {$d_\infty$ (Supremum)};
    
    \draw[eMediumBlue, very thick] (0,0) circle (1cm);
    \node[eMediumBlue] at (1.2, 0.5) {$d_2$ (Standard)};
    
    \draw[eDarkRed, very thick, dotted] (1,0) -- (0,1) -- (-1,0) -- (0,-1) -- cycle;
    \node[eDarkRed] at (0.7, -0.7) {$d_1$ (Manhattan)};
\end{tikzpicture}
\end{center}
\end{example*}

\begin{example*}[Continuous Functions]
Let $X = C^0([0,1], \R)$ denote the set of all continuous functions mapping the interval $[0,1]$ to $\R$. We can equip this function space with the \textcolor{eDarkGreen}{\qt{Supremum Metric}}:
\[ d(f,g) := \sup_{x \in [0,1]} \abs{f(x) - g(x)} \]
\textcolor{eGray}{Intuition: This metric finds the single point $x$ where the two graphs are furthest apart, and defines that maximum vertical gap as the distance between the two functions.}
\end{example*}

\begin{example*}[The n-Sphere]
Let $X = S^2 := \set{ (x_1, x_2, x_3) \in \R^3 : x_1^2 + x_2^2 + x_3^2 = 1 }$ (the 2-dimensional sphere). 
The closest path between two points on the sphere is an arc of a great circle (a geodesic). This defines a unique metric on the surface, similar to how straight lines define a metric on $\R^n$.

\begin{center}
\begin{tikzpicture}[scale=2]
    \draw[thick, eBlack] (0,0) circle (1cm);
    \draw[dashed, eBlack] (1,0) arc (0:180:1cm and 0.3cm);
    \draw[eBlack] (1,0) arc (0:-180:1cm and 0.3cm);
    
    \coordinate (X) at (-0.7, 0.4);
    \coordinate (Y) at (0.6, 0.6);
    
    \fill[eSaddleBrown] (X) circle (1.5pt) node[left] {$x$};
    \fill[eSaddleBrown] (Y) circle (1.5pt) node[right] {$y$};
    
    \draw[eDarkRed, dashed, thick] (X) -- node[below=0.2cm, sloped] {\small Euclidean distance} (Y);
    \draw[eMediumBlue, very thick] (X) to[bend left=20] node[above, sloped] {\small Distance on sphere} (Y);
\end{tikzpicture}
\end{center}
\geminiwarning{Terminology check: A sphere embedded in 3D space is actually a 2-dimensional surface, which is why topologists denote it as $S^2$. The script loosely refers to it as the \qt{n-sphere} here based on the ambient space.}
\end{example*}

\vspace{0.5cm}

\begin{exercise*}
Show that the supremum metric from the continuous functions example is indeed a metric.
\end{exercise*}

\begin{solution}
Symmetry and definiteness follow directly from the properties of the absolute value $\abs{\cdot}$. For the triangle inequality, let's consider three functions $f, g, h \in C^0([0,1])$ and define two helpful auxiliary functions:
\[ p := f - h, \quad q := h - g \]
For any arbitrary but fixed $x \in [0,1]$, the value of the function at that specific point can never exceed the supremum of the entire function:
\[ \abs{p(x)} \le \sup \abs{p} \quad \text{and} \quad \abs{q(x)} \le \sup \abs{q} \]
If we now look at the distance between $f$ and $g$ at the point $x$, we can use the "add zero" trick and apply our inequalities:
\begin{equation}
    \begin{aligned}
        \abs{f(x) - g(x)} &= \abs{p(x) + q(x)} \\
        &\le \abs{p(x)} + \abs{q(x)} \quad \text{\textcolor{eGray}{(Standard triangle inequality in $\R$)}} \\
        &\le \sup \abs{p} + \sup \abs{q} \\
        &= \sup \abs{f - h} + \sup \abs{h - g}
    \end{aligned}
\end{equation}
Because this upper bound holds for \textcolor{eSaddleBrown}{every single} $x \in [0,1]$, it must inherently also hold for the supremum taken over all $x$! Thus:
\[ d(f,g) \le d(f,h) + d(h,g) \]
This proves the triangle inequality.
\end{solution}

\hrulefill

\section{Open and Closed Sets}

Let $(X, d)$ be a metric space, let $x \in X$, and let $r \in \R_{>0}$. First, we define the \textcolor{eSaddleBrown}{\qt{open ball}} of radius $r$:
\[ B_r(x) := \set{ y \in X : d(y,x) < r } \]
\textcolor{eGray}{Note: Pay attention to the strict inequality ($<$). The boundary shell is not part of the ball!}

\begin{definition*}
\begin{itemize}
    \item A subset $U \subseteq X$ is \textcolor{eDarkGreen}{\textbf{open}} iff for \emph{every} $y \in U$, we can find an $\eps > 0$ such that the entire ball $B_\eps(y)$ lies completely inside $U$. \textcolor{eGray}{(You can never accidentally step on the boundary.)}
    \item A subset $A \subseteq X$ is \textcolor{eDarkRed}{\textbf{closed}} iff its complement $X \setminus A$ is open.
\end{itemize}
\end{definition*}

\begin{center}
\begin{tikzpicture}[scale=1.2, very thick]
    \draw[eMediumBlue, fill=eMediumBlue!10, dashed] (0,0) .. controls (1,2) and (3,1) .. (4,0) .. controls (3,-2) and (1,-1) .. (0,0);
    \node[eMediumBlue] at (3.5, 1) {\Large $U$};
    
    \coordinate (Y) at (1.5, 0);
    \draw[eDarkGreen, fill=eDarkGreen!20, dashed] (Y) circle (0.6cm);
    \fill[eSaddleBrown] (Y) circle (2pt) node[below] {$y$};
    \draw[eDarkGreen, ->] (Y) -- ++(45:0.6cm) node[midway, above left] {$\eps$};
    \node[eDarkGreen] at (2.5, -0.5) {$B_\eps(y) \subseteq U$};
\end{tikzpicture}
\end{center}

\vspace{0.5cm}

\begin{exercise*}
Let $(X,d)$ be a metric space. Show that the entire space $X$ and the empty set $\emptyset$ are both open and closed.
\end{exercise*}

\begin{solution}
First, let's look at openness:
\begin{itemize}
    \item $X$ is open because for any $x \in X$ and any radius $r > 0$, the ball $B_r(x)$ is by definition a subset of $X$. Thus, $B_r(x) \subseteq X$ trivially holds.
    \item $\emptyset$ contains no points. Therefore, the condition \qt{for every point $y \in \emptyset$ there exists...} is a vacuously true statement. Thus, the empty set is open.
\end{itemize}
Now for closedness:
Since $X$ and $\emptyset$ are each other's complements in $X$ (i.e., $X \setminus X = \emptyset$ and $X \setminus \emptyset = X$), and we just proved both are open, they must both inherently be closed as well!
\end{solution}

\hrulefill

\section{Sequences and Convergence}

A sequence $(x_n)_{n=0}^\infty$ in a metric space $(X, d)$ is formally a function from $\N \to X$. We say that $(x_n)$ \textbf{converges} to a limit $x \in X$ if:
\[ \forall \eps > 0, \exists N \in \N \text{ such that } d(x_n, x) < \eps \quad \forall n \ge N \]

\subsection*{The Big Connection Proposition}
Topological properties (openness/closedness) and the behavior of sequences go hand-in-hand:
\begin{enumerate}
    \item $U \subseteq X$ is \textbf{open} $\iff$ for every sequence in $X$ that has a limit in $U$, the sequence eventually lies entirely inside $U$.
    \item $A \subseteq X$ is \textbf{closed} $\iff$ \emph{every} convergent sequence that lies entirely in $A$ must necessarily have its limit in $A$ as well. \textcolor{eGray}{(The sequence cannot \qt{break out} of the set.)}
\end{enumerate}

\subsection*{\warning{Watch out! A Common Misconception!}}
Let's look at the standard Euclidean metric space $\R^2$. We know the square $[0,1]^2$ is closed and $(0,1)^2$ is open. (The rectangle $[0,1] \times (0,1)$ is neither open nor closed in $\R^2$).

\textcolor{eDarkRed}{But...} what if our \textit{entire universe} (our metric space $X$) is defined strictly as $X := [0,1]^2$? 
In the subspace $X = [0,1]^2$, the set $[0,1]^2$ is suddenly \emph{open}! 

\textcolor{eDarkCyan}{Why?} Because the full space $X$ is always open within itself. If you take a point exactly on the boundary, its open $\eps$-ball consists only of points \textcolor{eSaddleBrown}{that actually exist in $X$}. The ball does not bleed out into $\R^2$ because $\R^2$ simply does not exist for this specific metric space!

\begin{center}
\begin{tikzpicture}[scale=2, very thick]
    \draw[->, eBlack] (-0.5, 0) -- (1.5, 0) node[right] {$x_1$};
    \draw[->, eBlack] (0, -0.5) -- (0, 1.5) node[above] {$x_2$};
    
    \draw[eMediumBlue, fill=eMediumBlue!10] (0,0) rectangle (1,1);
    \node[eMediumBlue] at (0.5, 0.5) {$X = [0,1]^2$};
    
    \coordinate (P) at (1, 0.3);
    \fill[eDarkRed] (P) circle (1.5pt) node[right] {$x$};
    \draw[eDarkRed, dashed] (P) circle (0.3cm);
    
    \begin{scope}
        \clip (0,0) rectangle (1,1);
        \fill[eDarkRed!40] (P) circle (0.3cm);
    \end{scope}
    
    \node[eDarkRed, align=left, anchor=west] at (1.5, 0.5) {The $\eps$-ball in $X$ is \\ strictly the shaded region! \\ $B_\eps(x) \subseteq X$ is perfectly valid.};
\end{tikzpicture}
\end{center}

\begin{exercise*}
\begin{enumerate}
    \item Write down the definition of convergence for a sequence of functions $(f_n)_{n=0}^\infty$ to $f$ in the metric space $X$ of bounded functions $[0,1] \to \R$ equipped with the supremum metric.
    \item You know this resulting expression from Analysis 1. What kind of convergence is this called?
    \item Use the fact that the space of continuous functions $C^0([0,1]) \subseteq X$ is closed, and combine it with the \qt{sequentially closed} proposition (part b) to conclude a massive result from Analysis 1.
\end{enumerate}
\end{exercise*}

\begin{solution}
\begin{enumerate}
    \item $\forall \eps > 0 \, \exists N \in \N$ such that $\sup_{x \in [0,1]} \abs{f_n(x) - f(x)} < \eps \quad \forall n \ge N$.
    \item This is exactly the definition of \textcolor{eDarkGreen}{\textbf{Uniform Convergence}}!
    \item The fact that $C^0([0,1])$ is closed in the space of bounded functions means (according to our proposition above): If a sequence of continuous functions converges uniformly to a limit function $f$, then this limit function $f$ cannot leave the space of continuous functions. Ergo: \textbf{The uniform limit of continuous functions is always continuous!}
    
    \textcolor{eSaddleBrown}{Proof of closedness (The $\eps/3$-trick):} 
    We show that the complement is open. Let $f \in X \setminus C^0([0,1])$. This means $f$ is bounded but has at least one discontinuity at a point $x_0$. That implies there is a real jump $\eps > 0$, such that no matter how close ($\delta$) we get to $x_0$, we can always find a $y$ where $\abs{f(x_0) - f(y)} \ge \eps$.
    
    Now we consider an open ball around this discontinuous function $f$, but we give it a tiny radius of $\eps/3$:
    \[ B_{\eps/3}(f) = \set{g \in X : \sup_{x \in [0,1]} \abs{f(x) - g(x)} < \frac{\eps}{3}} \]
    
    \begin{center}
    \begin{tikzpicture}[scale=1.5, very thick]
        \draw[->, thick, eBlack] (-0.5, 0) -- (4, 0) node[right] {$x$};
        \draw[->, thick, eBlack] (0, -0.5) -- (0, 3) node[above] {$y$};
        \node[below left, eBlack] at (0,0) {$0$};
        \draw[thick, eBlack] (3.5, 0.1) -- (3.5, -0.1) node[below] {$1$};
        
        \draw[eMediumBlue] (0, 0.5) to[out=0, in=180] (1.5, 0.2);
        \draw[eMediumBlue] (1.5, 1.5) to[out=0, in=180] (3.5, 2.5);
        
        \fill[eMediumBlue] (1.5, 1.5) circle (1.5pt);
        \fill[eMediumBlue, fill=white, draw=eMediumBlue, thick] (1.5, 0.2) circle (1.5pt);
        
        \draw[<->, eSaddleBrown] (1.5, 0.3) -- (1.5, 1.4) node[midway, right] {Jump $\eps$};
        
        \draw[eDarkRed, dashed] (0, 0.8) to[out=0, in=180] (1.5, 0.5);
        \draw[eDarkRed, dashed] (0, 0.2) to[out=0, in=180] (1.5, -0.1);
        \draw[eDarkRed, dashed] (1.5, 1.8) to[out=0, in=180] (3.5, 2.8);
        \draw[eDarkRed, dashed] (1.5, 1.2) to[out=0, in=180] (3.5, 2.2);
        
        \draw[<->, eDarkRed] (2.5, 2.05) -- (2.5, 2.35) node[midway, right] {$\frac{\eps}{3}$};
        \node[eMediumBlue] at (3.7, 2.5) {$f$};
    \end{tikzpicture}
    \end{center}
    
    This ball (the \qt{$\eps/3$-tube} around $f$) contains exclusively functions that are forced to replicate this jump. No function $g$ trapped inside this tube can possibly be continuous! Therefore, the entire ball $B_{\eps/3}(f)$ is contained within the complement $X \setminus C^0([0,1])$. Since the complement is open, $C^0([0,1])$ is closed!
\end{enumerate}
\end{solution}

\hrulefill

\section*{Friday}

\begin{warmup}
What is the shape of the open unit ball $B_1(0) \subseteq \R^2$ with:
\begin{itemize}
    \item The Supremum metric?
    \item The Manhattan metric?
\end{itemize}
\end{warmup}

\begin{answer*}\hfill\\ 
\begin{center}
\begin{tikzpicture}[scale=1.5, very thick]
    \node[eDarkCyan, font=\bfseries] at (0, 1.5) {Supremum ($d_\infty$):};
    \draw[->, thick, eBlack] (-1.5,0) -- (1.5,0) node[right] {$x_1$};
    \draw[->, thick, eBlack] (0,-1.5) -- (0,1.5) node[above] {$x_2$};
    
    \draw[eDarkGreen, dashed] (-1,-1) rectangle (1,1);
    \fill[eSaddleBrown] (0,1) circle (1.5pt) node[above right] {$(0,1)$};
    \fill[eSaddleBrown] (1,0) circle (1.5pt) node[above right] {$(1,0)$};

    \begin{scope}[xshift=4cm]
        \node[eDarkCyan, font=\bfseries] at (0, 1.5) {Manhattan ($d_1$):};
        \draw[->, thick, eBlack] (-1.5,0) -- (1.5,0) node[right] {$x_1$};
        \draw[->, thick, eBlack] (0,-1.5) -- (0,1.5) node[above] {$x_2$};
        
        \draw[eDarkRed, dashed] (1,0) -- (0,1) -- (-1,0) -- (0,-1) -- cycle;
        \fill[eSaddleBrown] (0,1) circle (1.5pt) node[above right] {$(0,1)$};
        \fill[eSaddleBrown] (1,0) circle (1.5pt) node[above right] {$(1,0)$};
    \end{scope}
\end{tikzpicture}
\end{center}
\end{answer*}

\vspace{0.5cm}

\begin{definition*}[Continuous functions]
For two metric spaces $(X, d_X)$ and $(Y, d_Y)$, a function $f: X \to Y$ is called \emph{continuous} if one of the following, completely equivalent, conditions holds:

\begin{enumerate}
    \item \textbf{$(\eps-\delta)$-criterion:} $\forall x_0 \in X, \forall \eps > 0 \, \exists \delta > 0$ such that $\forall x \in X$:
    \[ d_X(x_0, x) < \delta \implies d_Y(f(x_0), f(x)) < \eps \]
    \textcolor{eGray}{(If you are close enough to $x_0$ in the starting space, you are guaranteed to land close to $f(x_0)$ in the target space.)}
    
    \item \textbf{Sequential continuity:} For any sequence $(x_n)_{n=0}^\infty$ in $X$ with a limit $x \in X$, it holds that:
    \[ \lim_{n \to \infty} f(x_n) = f(x) \]
    \textcolor{eGray}{(Continuous functions allow you to pull the limit inside the argument.)}
    
    \item \textbf{Topological continuity:} For all open sets $U \subseteq Y$, the preimage $f^{-1}(U) \subseteq X$ is \emph{open}.
    \textcolor{eGray}{(Continuous functions preserve the topological structure when mapping backwards.)}
\end{enumerate}
\end{definition*}

\begin{example*}
Let $f: \R \to \R$, $x \mapsto x^2$ and let $(a,b) \subseteq \R$ be an interval. The preimage of this interval depends heavily on $a$ and $b$:
\begin{itemize}
    \item For $a < 0 < b$: $f^{-1}((a,b)) = (-\sqrt{b}, \sqrt{b})$. \textcolor{eGray}{(This is an open interval!)}
    \item For $a < b < 0$: $f^{-1}((a,b)) = \emptyset$. \textcolor{eGray}{(The empty set is, as we proved, open!)}
    \item For $0 < a < b$: $f^{-1}((a,b)) = (-\sqrt{b}, -\sqrt{a}) \cup (\sqrt{a}, \sqrt{b})$. \textcolor{eGray}{(A union of open intervals is open!)}
\end{itemize}
\textcolor{eDarkCyan}{Why is it enough to check only these simple basis intervals?} \\
\textcolor{eGray}{Hint: The preimage perfectly respects unions: $f^{-1}(U_1 \cup U_2) = f^{-1}(U_1) \cup f^{-1}(U_2)$. Any open set on $\R$ can be constructed by patching together such intervals.}
\end{example*}

\vspace{0.5cm}

\begin{question*}
Consider functions mapping between two metric spaces $(X, d_X)$ and $(Y, d_Y)$.
\begin{enumerate}
    \item What happens if $d_Y$ is the discrete metric? Is a standard function like $f: (\R, d_{std}) \to (\R, d_{discrete})$ continuous?
    \item And conversely: What if the starting space $d_X$ is discrete, and the target space $d_Y$ is entirely arbitrary?
\end{enumerate}
\end{question*}

\begin{answer*}
\begin{enumerate}
    \item \textcolor{eDarkRed}{No, almost never.} Take the identity function $f(z) = z$ as an example. In the discrete space $Y$, even a single point like $\set{0}$ is an open set! Let's look at its preimage: $\text{Id}^{-1}(\set{0}) = \set{0}$. But a single point is definitely \textbf{not} open in the starting space with the standard metric $d_{std}$. The preimage of an open set is not open here, so the function is not continuous!
    
    \item \textcolor{eDarkGreen}{Yes, always!} In a discrete space, every single subset is open. This means that no matter what open set $U \subseteq Y$ you pick, its preimage $f^{-1}(U)$ is guaranteed to be open in $X$ (because \emph{everything} is open there). Any function mapping out of a discrete space is automatically continuous.
\end{enumerate}
\end{answer*}

\hrulefill

\section{Completeness and Compactness}

\begin{definition*}[Cauchy Sequences \& Complete Metric Spaces]
Let $(X, d)$ be a metric space, and $(x_n)_{n=0}^\infty$ a sequence in $X$. We call $(x_n)$ a \textcolor{eDarkGreen}{\textbf{Cauchy sequence}} if its terms get arbitrarily close to one another beyond a certain point:
\[ \forall \eps > 0 \, \exists N \in \N \text{ such that } d(x_n, x_m) < \eps \quad \forall n, m \ge N \]

A metric space $(X, d)$ is called a \textcolor{eDarkGreen}{\textbf{complete metric space}} if every Cauchy sequence in $X$ actually converges to a real limit within $X$.
\end{definition*}

\gemininote{Intuition for completeness: A Cauchy sequence desperately "wants" to converge. A space is incomplete if it sabotages this—for instance, if the space has \qt{holes} (like $\mathbb{Q}$) or is missing its boundary, meaning the sequence steers towards a limit that isn't actually there.}

\begin{example*}
\begin{itemize}
    \item The Euclidean spaces $(\R^n, \langle \cdot, \cdot \rangle_{\text{eucl}})$ are complete.
    \item Any non-closed subset $U \subseteq \R^n$ and the rational numbers $\mathbb{Q}$ are \emph{not} complete.
    \item The space of continuous functions $(C^0([0,1]), \sup)$ is complete!
\end{itemize}
\end{example*}

\begin{definition*}[Compactness]
A metric space $(X, d)$ is compact if every wild, infinite open cover can be reduced to a \emph{finite} subcover.
\end{definition*}

\begin{theorem*}
A metric space $(X, d)$ is compact if and only if \emph{every} sequence in $X$ has a subsequence that converges to a point in $X$.
\end{theorem*}

\begin{center}
\begin{tikzpicture}[scale=1.2, very thick]
    \draw[eMediumBlue, fill=eMediumBlue!5] (0.5,0.5) .. controls (-0.5,1.5) and (1,2.8) .. (2.5,2.5) .. controls (4.5,2.2) and (4,0) .. (2,0) .. controls (1,-0.5) and (0.8,0) .. cycle;
    \node[eMediumBlue] at (4.2, 2.5) {\Large $X$};
    
    \draw[eDarkRed, dashed, fill=eDarkRed!10, fill opacity=0.3] (1.2, 0.4) circle (1.4cm);
    \node[eDarkRed] at (1.2, -0.6) {$U_1$};
    
    \draw[eSaddleBrown, dashed, fill=eSaddleBrown!10, fill opacity=0.3] (2.8, 1.2) circle (1.6cm);
    \node[eSaddleBrown] at (3.8, 0.8) {$U_2$};
    
    \draw[eMediumOrchid, dashed, fill=eMediumOrchid!10, fill opacity=0.3] (1.0, 1.8) circle (1.3cm);
    \node[eMediumOrchid] at (0.0, 2.6) {$U_3$};
    
    \node[eDarkGray] at (2, -1.5) {$U_1 \cup U_2 \cup U_3 = X$};
\end{tikzpicture}
\end{center}

\textcolor{eGray}{Note: A subset $K$ of a metric space $(X, d)$ is compact if the subspace $(K, d)$ is compact. This is strictly equivalent to the following prominent definition from the official script:}

\vspace{0.5cm}

% --- MAGIC COUNTER TRICK FOR DEF 9.63 ---
\newcounter{tempSection}
\setcounter{tempSection}{\value{section}}
\newcounter{tempTheorem}
\setcounter{tempTheorem}{\value{theorem}}

\setcounter{section}{9}
\setcounter{theorem}{62} 

\begin{definition}[Compactness]
Let $(X, d)$ be a metric space. A subset $K \subseteq X$ is called compact if one (and therefore both) of the following equivalent conditions hold:
\begin{enumerate}
    \item $K$ is \textcolor{eMediumBlue}{\qt{sequentially compact}}: every sequence $(x_n)_{n \in \N}$ in $K$ has a subsequence that is convergent in $K$.
    \item $K$ is \textcolor{eMediumBlue}{\qt{topologically compact}}: every family of open sets $\set{U_i}_{i \in I}$ that cover $K$ has a \emph{finite} subcover.
\end{enumerate}
\end{definition}

% --- RESTORE COUNTERS ---
\setcounter{section}{\value{tempSection}}
\setcounter{theorem}{\value{tempTheorem}}

\begin{center}
\begin{tikzpicture}[scale=1.2, thick]
    \draw[eDarkGray, dashed] (-3, -2) rectangle (5, 3);
    \node[eDarkGray] at (4.5, 2.5) {$X$};

    \draw[eDarkRed, fill=eDarkRed!20, rounded corners=15pt] (0,0) -- (1,2) -- (3,1.5) -- (2,-1) -- cycle;
    \node[eDarkRed] at (1.5, 0.5) {\Large $K$};

    \draw[eMediumBlue, dashed, fill=eMediumBlue!10, fill opacity=0.5, very thick] (0.2, 0.8) circle (1.2cm);
    \draw[eMediumBlue, dashed, fill=eMediumBlue!10, fill opacity=0.5, very thick] (2.2, 0.5) circle (1.3cm);
    \draw[eMediumBlue, dashed, fill=eMediumBlue!10, fill opacity=0.5, very thick] (1.2, -0.2) circle (1cm);
    
    \node[eMediumBlue] at (-0.5, 1.5) {$U_{\alpha_1}$};
    \node[eMediumBlue] at (3.5, 1.2) {$U_{\alpha_2}$};
    \node[eMediumBlue] at (1.2, -1.4) {$U_{\alpha_3}$};
\end{tikzpicture}
\end{center}
\textcolor{eGray}{Important: The sets $U_i \subseteq X$ in the cover are always \emph{open sets} with respect to the ambient space $(X,d)$.}

\vspace{0.5cm}

\begin{exercises}
\begin{enumerate}
    \item Show that every finite subset $\set{x_1, \dots, x_n} \subseteq X$ is compact.
    \item If a set $K$ is compact, is it automatically complete?
\end{enumerate}
\end{exercises}

\begin{solution}\hfill\\ 
\textbf{a)} Let's use topological compactness: Let $\set{U_\alpha}_{\alpha \in I}$ be an arbitrary open cover of $\set{x_1, \dots, x_n}$. Since these sets must cover all our points, we can simply pick \emph{one} set $U_{\alpha_i}$ from the massive pool for each point $x_i$ (from $i=1$ to $n$) such that $x_i \in U_{\alpha_i}$. 
The collection $\set{U_{\alpha_1}, \dots, U_{\alpha_n}}$ is then a subcover. Because it consists of only $n$ (a finite number!) sets, we have found a finite subcover. Thus, the finite set is compact.

\vspace{0.3cm}
\textbf{b)} \textcolor{eDarkGreen}{Yes, absolutely.} Sequential compactness is perfect for this: Let $(x_n)_{n=0}^\infty$ be an arbitrary Cauchy sequence in $X$. Because $K$ is (sequentially) compact, this sequence \emph{must} have a convergent subsequence:
\[ \lim_{i \to \infty} x_{n_i} = x \in X \]
Now a fundamental trick of Cauchy sequences comes into play: If a Cauchy sequence agrees to converge even for a fraction of its terms (the subsequence), it forcibly drags the rest of the sequence along with it!

Formally: Fix an $\eps > 0$. Because it is a Cauchy sequence, all terms move within $\eps$-distance of each other beyond some index $N_1$:
\[ d(x_n, x_m) < \eps \quad \forall n, m > N_1 \]
Because the subsequence converges to $x$, its terms move within $\eps$-distance of $x$ beyond some index $N_2$:
\[ d(x, x_{n_i}) < \eps \quad \forall i > N_2 \]
If we go far enough into the sequence (choosing $N := \max\set{N_1, N_2}$), we can estimate the distance of any arbitrary term $x_k$ (with $k > N$) to the limit $x$. We just cleverly insert an element $x_{n_k}$ of the convergent subsequence in between:
\[ d(x, x_k) \le d(x, x_{n_k}) + d(x_{n_k}, x_k) < 2\eps \]
\textcolor{eGray}{(Where of course $n_k \ge k > N$).} Since we can make the distance $2\eps$ arbitrarily small, the entire Cauchy sequence indeed converges to $x$. Thus, every compact space is always complete!
\end{solution}